\origpage{109}{GRAPH MINORS. XIII}
\hspace*{12pt}(2) If we obtain a branch-decomposition of $G_i \setminus Z$ of width $\le \omega$, we convert it to a branch-decomposition of $G_i$ of width $\le \fix{\max(}\omega\fix{, 2)} + \xi$\fixnote{fix:10.5-step2}{The same correction as in Algorithm (10.1), for the same reason: if $G_i \setminus Z$ has at most one edge, its branch-width is 0 by convention, and adding the $\xi$ vertices of $Z$ can give branch-width 2.}, compute the $\delta$-folio of $G_i$ relative to $Z$ by (4.1), and stop.

\hspace*{12pt}(3) If we obtain a model of $K_p$ in $G_i \setminus Z$ and hence in $G_i$, we apply (6.2) to find a vertex $v$ of $G_i$ irrelevant to the $\delta$-folio of $G_i$ relative to $Z$; we set $G_{i+1} = G_i \setminus v$ and return for the next iteration.

\hspace*{12pt}(4) Otherwise, we obtain $X$, $H$, $\mathcal{A}$ as in Hypothesis C (with $G$ replaced by $G_i \setminus Z$), where $|X| \le \binom{p}{2} - 1$ and $H$ has height $h$. Put $X' = X \cup Z$; then $G_i$, $X'$, $H$, $\mathcal{A}$ also satisfy Hypothesis C. We apply (10.4) to $G_i$, $X'$, $H$, $\mathcal{A}$, to compute an irrelevant vertex $v$ for the $\delta$-folio of $G_i$ relative to $X'$. Since $Z \subseteq X'$, $v$ is also irrelevant for the $\delta$-folio of $G_i$ relative to $Z$, by (2.1). We set $G_{i+1} = G_i \setminus v$, and return for the next iteration.

The application of (6.5) in step (1) takes time $O(v^3)$. In the recursion, each iteration takes time $O(\varepsilon(G')^2)$. Since there are $\le v(G)$ iterations, and $\varepsilon(G') \le O(v(G))$, the algorithm takes time $O(v^3)$. $\square$

We would like to say that not only does there exist an algorithm as in (10.5), but we can \textit{construct} it, and that depends on being able to construct the various constants $p, h, \omega$ used in the course of the proof, for given numerical values of $\xi$ and $\delta$. Now $p$ is easy, but $h$ and $\omega$ are defined by (10.3) and (9.8), and these are formulated as purely existential statements, containing no method for calculating $h$ and $\omega$. (``For all $q$, $\delta \ge 0$ there exists $h \ge 0\dots$``) Nevertheless the proofs of (10.3) and (9.8) do yield procedures for calculating $h$ and $\omega$, except that (10.3) uses the function $f$ of (10.2), and (9.8) uses $h'$ of (9.5) and another $\omega$, from (4.2). The sequences of definitions can be traced further and further back, through other papers in this series, until they terminate; and all the proofs used to define these constants are constructive. Even, say, Theorem (3.5) of [16], which asserts that ``For every matching $G$ in $\Sigma$ and every homoplasty class $\mathcal{D}$ there exists $H \in \mathcal{D}$ with $U(G) \cap U(H)$ finite'' can be proved constructively, in the sense that it is easy to give an algorithm which, given explicitly some $G$ and $\mathcal{D}$, will compute an appropriate $H$. Thus, we do indeed have an algorithm which, with input numbers $\xi$ and $\delta \ge 0$, finds the algorithm of (10.5); and also the constant $C$ in the bound on its running time can be constructed.

\sectionhead{References}

1. S. Arnborg, D. G. Corneil, and A. Proskurowski, Complexity of finding embeddings in $k$-trees, \textit{SIAM J. Algebraic Discrete Methods} \textbf{8} (1987), 277-284.

2. S. Arnborg and A. Proskurowski, Linear time algorithms for NP-hard problems restricted to partial $k$-trees, \textit{Discrete Applied Math.} \textbf{23} (1989), 11-24.
